\documentclass[11pt,a4paper]{article}

\usepackage[a4paper,margin=2.5cm]{geometry}
\usepackage{fontspec}
\usepackage{titlesec}
\usepackage{longtable}
\usepackage{booktabs}
\usepackage{array}
\usepackage{ragged2e}
\usepackage{xcolor}
\usepackage{hyperref}
\usepackage{enumitem}
\usepackage{fancyhdr}

\setmainfont{Sylfaen}
\newfontfamily\latinfont{Georgia}

\definecolor{brandgreen}{HTML}{1F6F5C}
\definecolor{rowgray}{HTML}{F8F9FA}

\hypersetup{
  colorlinks=true,
  linkcolor=brandgreen,
  urlcolor=brandgreen,
  pdftitle={მოთხოვნათა კატალოგი -- LLM-Driven Climate-Health ETL Platform},
  pdfauthor={სერგი კონიაშვილი}
}

\titleformat{\section}{\large\bfseries\color{brandgreen}}{\thesection}{0.5em}{}
\titleformat{\subsection}{\normalsize\bfseries}{\thesubsection}{0.5em}{}

\setlength{\headheight}{14pt}
\pagestyle{fancy}
\fancyhf{}
\fancyhead[L]{LLM-Driven Climate-Health ETL Platform}
\fancyhead[R]{მოთხოვნათა კატალოგი}
\fancyfoot[C]{\thepage}

\setlength{\emergencystretch}{3em}
\setlength{\tabcolsep}{4pt}
\renewcommand{\arraystretch}{1.3}
\sloppy

\newcolumntype{L}[1]{>{\RaggedRight\arraybackslash}p{#1}}

\begin{document}

%=====================================================================
% სატიტულო გვერდი
%=====================================================================
\begin{titlepage}
\centering
\vspace*{2.5cm}

{\huge\bfseries\color{brandgreen} მოთხოვნათა კატალოგი\par}
\vspace{0.4cm}
{\Large (Anforderungskatalog / Requirements Catalog)\par}
\vspace{1.2cm}

{\large\textbf{პროექტი:}\par}
\vspace{0.2cm}
{\Large LLM-ზე დაფუძნებული კლიმატისა და ჯანმრთელობის მონაცემთა ინტეგრაციისა და ანალიზის პლატფორმა\par}
\vspace{0.3cm}
{\normalsize (LLM-Driven Climate-Health ETL Platform)\par}

\vspace{2.5cm}

\begin{tabular}{ll}
\textbf{ავტორი:} & სერგი კონიაშვილი \\
\textbf{სტატუსი:} & ბაკალავრიატის ნაშრომი (Bachelorarbeit) \\
\textbf{ვერსია:} & 2.1 (სრული ფუნქციონალი) \\
\textbf{თარიღი:} & 2026 წლის ოქტომბერი \\
\end{tabular}

\vfill

{\footnotesize
დოკუმენტი წარმოადგენს პროგრამული უზრუნველყოფის სრულ მოთხოვნათა კატალოგს,\\
რომელიც მოიცავს 23 ფუნქციონალურ (FR-01 -- FR-23) და 9 არაფუნქციონალურ (NFR-01 -- NFR-09) მოთხოვნას.\\
ყველა მოთხოვნა სრულად არის დანერგილი და შემოწმებული.
\par}

\vspace*{1.5cm}
\end{titlepage}

%=====================================================================
% სარჩევი
%=====================================================================
\tableofcontents
\newpage

%=====================================================================
% 1. შესავალი და მიზნები
%=====================================================================
\section{შესავალი და მიზნები}

მრავალი ინფექციური დაავადება (განსაკუთრებით ვექტორული და წყლით გადამცემი პათოლოგიები) მჭიდრო კორელაციაშია კლიმატურ ფაქტორებთან. მათი ბიოლოგიური გადამტანები (მაგალითად, {\latinfont\itshape Aedes aegypti} და {\latinfont\itshape Anopheles} სახეობის კოღოები) გამოირჩევიან მაღალი მგრძნობელობით ჰაერის ტემპერატურის, ნალექების ინტენსივობისა და ფარდობითი ტენიანობის მიმართ. მაღალი სიზუსტის მეტეოროლოგიური მონაცემები წარმოადგენს დაავადების გადაცემის ხელსაყრელობის უმნიშვნელოვანეს ადრეულ ინდიკატორს. მიუხედავად ამისა, ეპიდემიოლოგიურ მონაცემთა ნაკრებებთან კლიმატური მონაცემების შერწყმა მკვლევრებისა და საზოგადოებრივი ჯანდაცვის სპეციალისტებისთვის სერიოზულ სირთულეს ქმნის შემდეგი ბარიერების გამო:
\begin{itemize}[leftmargin=1.5em]
  \item ფრაგმენტული მონაცემთა წყაროები და ტექნიკური წვდომის შეზღუდვები,
  \item შეუსაბამო დროითი (დღიური კლიმატი vs.\ თვიური/წლიური შემთხვევები) და სივრცითი რეზოლუციები,
  \item რთული და არაერთგვაროვანი მონაცემთა ფორმატები,
  \item მოსახლეობის რაოდენობის (დემოგრაფიული ინდიკატორების) გათვალისწინების სირთულე ინციდენტობის გამოსათვლელად,
  \item კლიმატის ზემოქმედების დროითი დაყოვნება (ლაგ-ეფექტი), რომლის გაანგარიშებაც სპეციალურ სტატისტიკურ აპარატს მოითხოვს.
\end{itemize}

\textbf{პლატფორმის ძირითადი მიზანია} აღნიშნული ბარიერების სრული გადალახვა. სისტემა საშუალებას აძლევს მკვლევრებს, ჯანდაცვის ორგანოებსა და დაინტერესებულ პირებს (პროგრამირების ან მეტეოროლოგიის სიღრმისეული ცოდნის გარეშე), ბუნებრივ ენაზე მოთხოვნის შეყვანით ან ინტერაქტიული მართვის პანელით მიიღონ ანალიზისთვის გამზადებული, კლიმატთან ჰარმონიზებული ეპიდემიოლოგიური მონაცემები. დიდი ენობრივი მოდელი (LLM) მომხმარებლის მიზანს გარდაქმნის შესრულებად ETL საინტეგრაციო გეგმად, ახორციელებს ვერიფიცირებული სამეცნიერო წყაროების გამოკითხვას, დროითი და სივრცითი რეზოლუციების გათანაბრებას, ინციდენტობის გამოთვლას 100\,000 მოსახლეზე და ლაგ-კორელაციების სტატისტიკურ ანალიზს გასაგები განმარტებებით.

\textbf{ვალიდაციის ძირითადი სცენარი:} დენგეს ცხელება ტაილანდში (1990--2023), პლატფორმის სრული არქიტექტურული მზაობითა და გაფართოებით მსოფლიოს 39 ქვეყანასა და სხვა ინფექციურ დაავადებებზე (მალარია, ქოლერა, დასავლეთ ნილოსის ვირუსი და სხვ.).

%=====================================================================
% 2. ფუნქციონალური მოთხოვნები
%=====================================================================
\section{ფუნქციონალური მოთხოვნები (FR)}

\begin{longtable}{@{}>{\bfseries}p{1.6cm} L{7.4cm} L{3.6cm} L{2.4cm}@{}}
\toprule
ID & აღწერა & წარმომავლობა & სტატუსი \\
\midrule
\endhead
FR-01 & სისტემამ ბუნებრივ ენაზე უნდა მიიღოს მოთხოვნა (დაავადება, რეგიონი, კლიმატური ცვლადები, დროის პერიოდი) და LLM-ის მეშვეობით გარდაქმნას სტრუქტურირებულ საინტეგრაციო გეგმად. & პროექტის კონცეფცია & დანერგილია \\
FR-02 & მსოფლიო დაყოფილი უნდა იყოს ასარჩევ რეგიონებად; მომხმარებელს უნდა შეეძლოს შედეგების რეგიონების მიხედვით მიღება (ინტეგრირებულია 39 ქვეყანა). & ხელმძღვანელის მოთხოვნა & დანერგილია \\
FR-03 & სისტემამ თვალსაჩინოდ უნდა აღრიცხოს და გამოაჩინოს მილსადენის (ETL) ყოველი ეტაპი (გამოთხოვა, ვალიდაცია, ჰარმონიზაცია, დაკავშირება, სტატისტიკა). & ხელმძღვანელის მოთხოვნა & დანერგილია \\
FR-04 & კლიმატური მონაცემები ექსკლუზიურად ოფიციალური სამეცნიერო წყაროებიდან უნდა მოდიოდეს (Open-Meteo ERA5, NASA POWER, TMD). & ხელმძღვანელის მოთხოვნა & დანერგილია \\
FR-05 & მომხმარებელს უნდა ჰქონდეს შესაძლებლობა, აირჩიოს რამდენიმე ოფიციალურ სამეცნიერო მონაცემთა წყაროს შორის. & ხელმძღვანელის მოთხოვნა & დანერგილია \\
FR-06 & მომხმარებელს უნდა შეეძლოს საკუთარი გარე მონაცემთა წყაროს მითითება ვებ-ბმულის (URL) მეშვეობით (შემთხვევები, ამინდი, დემოგრაფია). & ხელმძღვანელის მოთხოვნა & დანერგილია \\
FR-07 & მომხმარებელს უნდა შეეძლოს საკუთარი მონაცემთა ნაკრების ატვირთვა CSV ფაილის სახით (შემთხვევები, ამინდი, დემოგრაფია). & ხელმძღვანელის მოთხოვნა & დანერგილია \\
FR-08 & მიღებული და გენერირებული მონაცემთა ნაკრებები მომხმარებელმა უნდა შეეძლოს ჩამოტვირთოს სტანდარტულ CSV და JSON ფორმატებში. & ხელმძღვანელის მოთხოვნა & დანერგილია \\
FR-09 & სისტემამ ყოველი გადამუშავების ეტაპი და კლიმატურ-ეპიდემიოლოგიური კავშირი უნდა ახსნას გასაგებ, ადამიანურ ენაზე. & პროექტის კონცეფცია & დანერგილია (LLM) \\
FR-10 & მონაცემთა გამოთხოვა და ვალიდაცია უნდა ხდებოდეს API-ს ლიმიტების დაცვითა და ასინქრონულად, რათა ინტერფეისი არ გაიყინოს. & მომხმარებლის მოთხოვნა & დანერგილია \\
FR-11 & სისტემამ უკვე დამუშავებული მონაცემები უნდა შეინახოს ქეშში და აჩვენოს ბოლო შემოწმების დრო (\texttt{last\_verified}). & მომხმარებლის მოთხოვნა & დანერგილია \\
FR-12 & არქიტექტურა უნდა იყოს მოდულარული, რათა შესაძლებელი იყოს ახალი დაავადებებისა და რეგიონების მარტივად დამატება. & პროექტის კონცეფცია & დანერგილია \\
FR-13 & სისტემას უნდა ჰქონდეს ინტერაქტიული Copilot ჩატ-ასისტენტი, რომელიც პასუხობს დარგობრივ კითხვებს, ავსებს ფორმის პარამეტრებს და პირდაპირ უშვებს ძიებას. & გაფართოება / UX & დანერგილია \\
FR-14 & სისტემას უნდა გააჩნდეს სამეცნიერო ლიტერატურის ცოდნის ბაზა (WHO, Lancet, IPCC, Mordecai et al.), შედარებითი ანალიზი და RAG-კონტექსტუალიზაცია. & დარგობრივი მოთხოვნა & დანერგილია \\
FR-15 & სისტემამ ატვირთული სამეცნიერო PDF ფაილებიდან ან ვებ-ბმულებიდან ხელოვნური ინტელექტით ავტომატურად უნდა ამოიღოს და დაასტრუქტურიროს მეტამონაცემები. & გაფართოება / AI & დანერგილია \\
FR-16 & სისტემამ უნდა უზრუნველყოს გარე სათვალთვალო წყაროების პირდაპირი ონლაინ ძიება (WHO GHO ინდიკატორები და HDX მონაცემები). & გაფართოება / მონაცემები & დანერგილია \\
FR-17 & სისტემამ უნდა მიაბას ოფიციალური დემოგრაფიული მონაცემები (მსოფლიო ბანკი, UN WPP, ატვირთვა) და გამოთვალოს ინციდენტობა 100\,000 მოსახლეზე. & დარგობრივი მოთხოვნა & დანერგილია \\
FR-18 & მომხმარებლის მიერ ატვირთული ფაილებისთვის სისტემამ უნდა აწარმოოს ტრანსფორმაციის უწყვეტი აუდიტის ჟურნალი (Audit Trail) და გასაგებად განმარტოს იგი. & ხარისხის მოთხოვნა & დანერგილია \\
FR-19 & სისტემამ უნდა გამოთვალოს დროითი დაყოვნების (Lags 0--3) ჯვარედინი კორელაციები პირსონის $r$-ითა და $p$-value სტატისტიკური სანდოობით. & ანალიტიკური მოთხოვნა & დანერგილია \\
FR-20 & დღიურ რეზოლუციაზე გადაყვანისას სისტემამ უნდა განახორციელოს მასის შემნახველი გლუვი დროითი ინტერპოლაცია და მომხმარებელს მისცეს მეთოდური რეკომენდაცია. & მეთოდოლოგიური მოთხოვნა & დანერგილია \\
FR-21 & სისტემას უნდა ჰქონდეს სრულად სამენოვანი ინტერფეისი (ქართული, გერმანული, ინგლისური) რეალურ დროში ენის მყისიერი გადართვით. & მომხმარებლის მოთხოვნა & დანერგილია \\
FR-22 & სისტემას უნდა შეეძლოს ბეჭდვისთვის მზა, მაღალი ხარისხის სამეცნიერო PDF ანალიტიკური ანგარიშის გენერირება (დიაგრამებით, ცხრილებითა და ციტატებით). & მომხმარებლის მოთხოვნა & დანერგილია \\
FR-23 & სისტემამ უნდა უზრუნველყოს ინტერაქტიული მსოფლიო რუკა ქვეყნებისა და კოორდინატების ვიზუალური, ინტუიციური არჩევისთვის. & მოხერხებულობის მოთხოვნა & დანერგილია \\
\bottomrule
\end{longtable}

%=====================================================================
% 3. არაფუნქციონალური მოთხოვნები
%=====================================================================
\section{არაფუნქციონალური მოთხოვნები (NFR)}

\begin{longtable}{@{}>{\bfseries}p{1.6cm} L{7.4cm} L{3.6cm} L{2.4cm}@{}}
\toprule
ID & აღწერა & წარმომავლობა & სტატუსი \\
\midrule
\endhead
NFR-01 & ლიმიტების დაცვა: სისტემა არ უნდა გასცდეს LLM API-ს დადგენილ კვოტებს (მოცურავე ფანჯრის მექანიზმი, მაქს. 15 მოთხოვნა/წთ). & საოპერაციო მოთხოვნა & დანერგილია \\
NFR-02 & სამეცნიერო ციტირებადობა: ყოველი პასუხი და ანგარიში ზუსტად უნდა უთითებდეს გამოყენებულ პირველად წყაროებს (კლიმატი, შემთხვევები, დემოგრაფია, ლიტერატურა). & სამეცნიერო სტანდარტები & დანერგილია \\
NFR-03 & პორტაბელურობა: სრული პლატფორმა (FastAPI ბექენდი და Flutter Web ფრონტენდი) გაშვებადი უნდა იყოს ერთიან Docker კონტეინერში. & DevOps / პორტაბელურობა & დანერგილია \\
NFR-04 & მონაცემთა უსაფრთხოება: API გასაღებები და კონფიდენციალური პარამეტრები არასდროს არ უნდა მოხვდეს კოდის საჯარო საცავში ან კლიენტის ბანდლში. & უსაფრთხოების სტანდარტები & დანერგილია (\texttt{.env}, \texttt{.gitignore}) \\
NFR-05 & აღქმადობა: გენერირებული ახსნა-განმარტებები გასაგები უნდა იყოს მომხმარებლისთვის ეპიდემიოლოგიაში ან მეტეოროლოგიაში სპეციალური განათლების გარეშე. & მომხმარებელზე ორიენტაცია & დანერგილია \\
NFR-06 & მხარდაჭერადობა: მონაცემთა მოპოვების, ჰარმონიზაციის, სტატისტიკური ანალიზისა და ვიზუალიზაციის მკაცრი დაყოფა დამოუკიდებელ მოდულებად. & პროგრამული ინჟინერია & დანერგილია \\
NFR-07 & RAG მოდელის სანდოობა: სისტემა იყენებს In-Context Grounding-ს ნეირონული წონების შეცვლის გარეშე, რაც გამორიცხავს ციტატების ჰალუცინაციას. & ხელოვნური ინტელექტის ეთიკა & დანერგილია \\
NFR-08 & მათემატიკური მასის შენახვა: დღიურ დონეზე ჩამოყვანის ინტერპოლაცია მკაცრად ინარჩუნებს ოფიციალური შეტყობინებების საერთო ჯამურ შემთხვევებს. & მონაცემთა მთლიანობა & დანერგილია \\
NFR-09 & ენობრივი თანმიმდევრულობა: ინტერფეისის ტექსტების სრული გამართულობა სამივე ენაზე (ქართული, გერმანული, ინგლისური) ენათა არევის გარეშე. & ენობრივი სტანდარტი (i18n) & დანერგილია \\
\bottomrule
\end{longtable}

%=====================================================================
% 4. არქიტექტურული მიმოხილვა
%=====================================================================
\section{არქიტექტურული მიმოხილვა}

პლატფორმა დაფუძნებულია სამდონიან, მოდულარულ არქიტექტურაზე:
\begin{enumerate}[label=(\arabic*)]\sloppy
  \item \textbf{პრეზენტაციის დონე (Flutter Web):} რეაქტიული, ერთგვერდიანი ვებ-აპლიკაცია სამენოვანი მხარდაჭერით, ინტეგრირებული Copilot ჩატის პანელით, ლიტერატურის კორპუსის მართვით, მონაცემთა ტრანსფორმაციის აუდიტის ჩვენებით, ინტერაქტიული მსოფლიო რუკითა და PDF ანგარიშის ექსპორტით.
  \item \textbf{სერვისებისა და ორკესტრაციის დონე (Python, FastAPI):} მოდულარული სერვისები მილსადენის მართვისთვის (\texttt{etl}), კლიმატური მონაცემების გამოთხოვისთვის (\texttt{climate}, \texttt{tmd}), ეპიდემიოლოგიური მონაცემების დამუშავებისთვის (\texttt{case\_data}, \texttt{who\_gho}, \texttt{hdx}), დემოგრაფიისთვის (\texttt{population}), ლაგ-სტატისტიკისთვის (\texttt{statistics}), ლიტერატურის ავტომატური ექსტრაქციისთვის (\texttt{literature\_extractor}) და LLM Copilot სერვისებისთვის (\texttt{llm}, \texttt{chat}).
  \item \textbf{ავტორიტეტული პირველადი წყაროების დონე:} კლიმატის რეანალიზის არქივები (Open-Meteo ERA5, NASA POWER, TMD), ეპიდემიოლოგიური მონაცემთა ბაზები (OpenDengue, WHO GHO, HDX), დემოგრაფიული ინდიკატორები (მსოფლიო ბანკი WDI, გაეროს მოსახლეობის დივიზიონი UN WPP) და Google Gemini LLM API.
\end{enumerate}

%=====================================================================
% 5. ვალიდაცია და ექსპერიმენტული შემოწმება
%=====================================================================
\section{ვალიდაცია და ექსპერიმენტული შემოწმება}

პლატფორმის მილსადენის ეფექტურობა ემპირიულად შემოწმდა ტაილანდში დენგეს ცხელების ისტორიულ მონაცემებზე (1990--2023 წლები). სისტემის არქიტექტურული მოქნილობა უზრუნველყოფს მის მყისიერ გამოყენებას მონაცემთა ბაზაში რეგისტრირებულ ყველა 39 ქვეყანაში, აგრეთვე მომხმარებლის მიერ ატვირთულ ნებისმიერ ახალ ლოკალურ თუ გლობალურ მონაცემთა ნაკრებზე.

\end{document}
